% Unit 037 English translation of chapter3.tex lines 946--1060.
% Source: Methods of Algebra, Volume 2, commit
% 9a5803ff77dd3257484cb177f851a73770a59dd3 (CC BY 4.0).
% stable-unit-id: o014.aljabr2.chapter3.abelian-category-complexes
% Coverage: Section 3.6 in full.

% segment-id: o014.li2.chapter3.abelian-category-complexes.h001
\section{Complexes on an Abelian Category}\label{sec:Abel-cplx}
\hypertarget{o014.aljabr2.chapter3.abelian-category-complexes}{ }

% segment-id: o014.li2.chapter3.abelian-category-complexes.p001
In this section, $\mathcal{A}$ is assumed to be an abelian category. As usual,
all statements about additivity in this section also extend to the case in
which $\mathcal{A}$ is a $\Bbbk$-linear category.

% segment-id: o014.li2.chapter3.abelian-category-complexes.p002
For a complex $X$ on an abelian category, we can discuss its cohomology
(Definition~\ref{def:cohomology}). Recall that $[n]$ not only shifts the
superscripts of a complex but also multiplies $d_X^\bullet$ by the sign
$(-1)^n$. This sign change does not alter the kernel or image of any $d_X$.
Thus
% segment-id: o014.li2.chapter3.abelian-category-complexes.q001
\[ \Hm^k\left( X[n] \right) = \Hm^{n+k} \left(X \right), \quad k, n \in \Z . \]

% segment-id: o014.li2.chapter3.abelian-category-complexes.p003
Proposition~\ref{prop:additive-cat-cplx} has already shown that
$\cate{C}(\mathcal{A})$ is an additive category. We now show that it is also
abelian.

% segment-id: o014.li2.chapter3.abelian-category-complexes.pr001
\begin{proposition}\label{prop:abelian-cat-cplx}
	The category $\cate{C}(\mathcal{A})$ is abelian. More precisely, for all
	$X,Y\in\Obj(\cate{C}(\mathcal{A}))$,
	$X \oplus Y = (X^n \oplus Y^n, (d_X^n, d_Y^n))_{n \in \Z}$, and for
	every morphism $f:X\to Y$ one may take
% segment-id: o014.li2.chapter3.abelian-category-complexes.q002
	\begin{gather*}
		\Ker(f) = \left( \Ker(f^n) \right)_{n \in \Z}, \quad
		\Coker(f) = \left( \Coker(f^n) \right)_{n \in \Z}, \\
		\Image(f) = \left( \Image(f^n) \right)_{n \in \Z}, \quad
		\Coim(f) = \left( \Coim(f^n) \right)_{n \in \Z}.
	\end{gather*}%
	\footnote{Translator's note (O014-C040): the source writes
	$\Ker(f)^n$, $\Coker(f)^n$, $\Image(f)^n$, and $\Coim(f)^n$ on the left
	but equates each with an entire family indexed by $n$ on the right. This
	translation removes the superscript $n$ on the left so that all four are
	identities of complexes; their degree-$n$ components remain exactly those
	shown on the right.}

	Suppose morphisms $f$ and $g$ in $\cate{C}(\mathcal{A})$ are composable
	and $gf=0$. Then
	$H:=\Hm\left[X\xrightarrow{f}Y\xrightarrow{g}Z\right]$ is the following
	complex:
% segment-id: o014.li2.chapter3.abelian-category-complexes.l001
	\begin{itemize}
% segment-id: o014.li2.chapter3.abelian-category-complexes.i001
		\item its degree-$n$ term is
		$H^n:=\Hm\left[X^n\xrightarrow{f^n}Y^n\xrightarrow{g^n}Z^n\right]$;
% segment-id: o014.li2.chapter3.abelian-category-complexes.i002
		\item the morphism $d_H^n:H^n\to H^{n+1}$ is determined by $d_X^n$,
		$d_Y^n$, and $d_Z^n$, together with the functoriality of
		$\Hm[\cdots]$ (Proposition~\ref{prop:homology-functoriality}).%
		\footnote{Translator's note (O014-C041): the source mentions only
		$d_X^n$ and $d_Y^n$. This translation adds $d_Z^n$, which is needed to
		give a morphism between the two three-term diagrams and hence to induce
		$d_H^n$.}
	\end{itemize}
	Consequently, $X\xrightarrow{f}Y\xrightarrow{g}Z$ is exact if and only if
	$X^n\xrightarrow{f^n}Y^n\xrightarrow{g^n}Z^n$ is exact for every
	$n\in\Z$.
\end{proposition}

% segment-id: o014.li2.chapter3.abelian-category-complexes.pf001
\begin{proof}
	Proposition~\ref{prop:additive-cat-cplx} explained how to use the forgetful
	functor from complexes to graded objects,
	$U:\cate{C}(\mathcal{A})\to\mathcal{A}^{\Z}$, to reduce the required
	$\varinjlim$ and $\varprojlim$ to degreewise constructions in
	$\mathcal{A}$. The point is that $U$ creates $\varinjlim$ and
	$\varprojlim$ (Lemma~\ref{prop:limits-cplx}). This gives degreewise
	constructions of $X\oplus Y$, $\Ker(f)$, $\Coker(f)$, and so forth.

	In particular, by the characterization \eqref{eqn:strict-morphism}, the
	canonical morphism $\Coim(f)\to\Image(f)$ is given by the morphisms
	$\Coim(f^n)\to\Image(f^n)$ in $\mathcal{A}$. A degreewise isomorphism is
	an isomorphism of complexes. Thus every morphism in
	$\cate{C}(\mathcal{A})$ is strict. Therefore
	$\cate{C}(\mathcal{A})$ is an abelian category.%
	\footnote{Translator's note (O014-C042): the closing sentence of the
	source names $\mathcal{A}$, although that category was already assumed
	abelian; the category just proved abelian is $\cate{C}(\mathcal{A})$.
	This translation corrects the category name.} The description of
	$\Hm[X\to Y\to Z]$ follows in the same way.
\end{proof}

% segment-id: o014.li2.chapter3.abelian-category-complexes.r001
\begin{remark}
	The same method proves that the category of double complexes
	$\cate{C}^2(\mathcal{A})$, and even the category of $k$-fold complexes
	$\cate{C}^k(\mathcal{A})$, remains abelian for $k\in\Z_{\geq1}$; see
	Remark~\ref{rem:k-fold-cplx}.
\end{remark}

% segment-id: o014.li2.chapter3.abelian-category-complexes.p004
Next consider a morphism $f:X\to Y$ in $\cate{C}(\mathcal{A})$. For every
$n\in\Z$, Proposition~\ref{prop:homology-functoriality} uniquely determines
$\Hm^n(f):\Hm^n(X)\to\Hm^n(Y)$ so that the following diagram commutes:
% segment-id: o014.li2.chapter3.abelian-category-complexes.g001
\begin{equation}\label{eqn:induced-morphism-cohomology}\begin{tikzcd}
	\Ker(d_X^n) \arrow[twoheadrightarrow, r] \arrow[d, "\text{induced by $f$}"'] & \Hm^n(X) \arrow[hookrightarrow, r] \arrow[d, "{\Hm^n(f)}"] & \Coker(d_X^{n-1}) \arrow[d, "\text{induced by $f$}"] \\
	\Ker(d_Y^n) \arrow[twoheadrightarrow, r] & \Hm^n(Y) \arrow[hookrightarrow, r] & \Coker(d_Y^{n-1}) .
\end{tikzcd}\end{equation}

% segment-id: o014.li2.chapter3.abelian-category-complexes.p005
For composable morphisms $X\xrightarrow{f}Y\xrightarrow{g}Z$, this
characterization immediately gives
$\Hm^n(gf)=\Hm^n(g)\Hm^n(f)$; moreover,
$\Hm^n(\identity_X)=\identity_{\Hm^n(X)}$.

% segment-id: o014.li2.chapter3.abelian-category-complexes.pr002
\begin{proposition}[Cohomology as a functor]\label{prop:H-functor}
	\index[sym1]{Hn}
	For every $n\in\Z$, the definition above gives an additive functor
	$\Hm^n:\cate{C}(\mathcal{A})\to\mathcal{A}$.
\end{proposition}

% segment-id: o014.li2.chapter3.abelian-category-complexes.pf002
\begin{proof}
	The properties $\Hm^n(gf)=\Hm^n(g)\Hm^n(f)$ and
	$\Hm^n(\identity)=\identity$ follow directly from the characterization
	\eqref{eqn:induced-morphism-cohomology}. In the same way one obtains
	$\Hm^n(f_1+f_2)=\Hm^n(f_1)+\Hm^n(f_2)$ and
	$\Hm^n(tf)=t\Hm^n(f)$ when $\mathcal{A}$ is $\Bbbk$-linear and
	$t\in\Bbbk$.
\end{proof}

% segment-id: o014.li2.chapter3.abelian-category-complexes.p006
A short exact sequence of complexes automatically induces a long exact
sequence in cohomology. This is the most elementary form of a long exact
sequence in homology theory.

% segment-id: o014.li2.chapter3.abelian-category-complexes.pr003
\begin{proposition}[A short exact sequence induces a long exact sequence]\label{prop:long-exact-sequence-ses}
	\index{long exact sequence}
	Let $0\to X\xrightarrow{f}Y\xrightarrow{g}Z\to0$ be a short exact
	sequence in $\cate{C}(\mathcal{A})$. Then Theorem~\ref{prop:snake-lemma}
	gives the following canonical exact sequence in $\mathcal{A}$:

% segment-id: o014.li2.chapter3.abelian-category-complexes.g002
	\begin{equation*}\begin{tikzcd}[row sep=large]
	\cdots \arrow[r] & \Hm^{n-1}(Y) \arrow[r, "{\Hm^{n-1}(g)}" inner sep=0.7em] \arrow[phantom, d, ""{coordinate, name=A}] & \Hm^{n-1}(Z) \arrow[dll, rounded corners, "{\delta^{n-1}}" description, to path={
		-- ([xshift=2ex]\tikztostart.east)
		|- (A) [near end]\tikztonodes
		-| ([xshift=-2ex]\tikztotarget.west)
		-- (\tikztotarget)}] \\
	\Hm^n(X) \arrow[r, "{\Hm^n(f)}"] & \Hm^n(Y) \arrow[r, "{\Hm^n(g)}"] \arrow[phantom, d, ""{coordinate, name=B}] & \Hm^n(Z) \arrow[dll, rounded corners, "{\delta^n}" description, to path={
		-- ([xshift=2ex]\tikztostart.east)
		|- (B) [near end]\tikztonodes
		-| ([xshift=-2ex]\tikztotarget.west)
		-- (\tikztotarget)}] \\
	\Hm^{n+1}(X) \arrow[r, "{\Hm^{n+1}(f)}"' inner sep=0.7em] & \Hm^{n+1}(Y) \arrow[r] & \cdots
	\end{tikzcd}\end{equation*}

	This sequence has the following functoriality. If there is a commutative
	diagram in $\cate{C}(\mathcal{A})$ whose rows are exact,
% segment-id: o014.li2.chapter3.abelian-category-complexes.g003
	\[\begin{tikzcd}[row sep=small]
		0 \arrow[r] & X \arrow[r] \arrow[d] & Y \arrow[r] \arrow[d] & Z \arrow[d] \arrow[r] & 0 \\
		0 \arrow[r] & \underline{X} \arrow[r] & \underline{Y} \arrow[r] & \underline{Z} \arrow[r] & 0
	\end{tikzcd}\]
	then one obtains a commutative diagram
% segment-id: o014.li2.chapter3.abelian-category-complexes.g004
	\[\begin{tikzcd}[column sep=small, row sep=small]
		\cdots \arrow[r] & \Hm^n(X) \arrow[d] \arrow[r] & \Hm^n(Y) \arrow[r] \arrow[d] & \Hm^n(Z) \arrow[r] \arrow[d] & \Hm^{n+1}(X) \arrow[r] \arrow[d] & \cdots \\
		\cdots \arrow[r] & \Hm^n(\underline{X}) \arrow[r] & \Hm^n(\underline{Y}) \arrow[r] & \Hm^n(\underline{Z}) \arrow[r] & \Hm^{n+1}(\underline{X}) \arrow[r] & \cdots .
	\end{tikzcd}\]
\end{proposition}

% segment-id: o014.li2.chapter3.abelian-category-complexes.pf003
\begin{proof}
	For every $n\in\Z$ there is an exact sequence
	$\Coker(d_X^{n-1})\to\Coker(d_Y^{n-1})\to\Coker(d_Z^{n-1})\to0$.
	Indeed, take the $\Coker$ part of the exact sequence in
	Theorem~\ref{prop:snake-lemma} applied to the diagram
% segment-id: o014.li2.chapter3.abelian-category-complexes.g005
	\[\begin{tikzcd}
		0 \arrow[r] & X^{n-1} \arrow[r] \arrow[d] & Y^{n-1} \arrow[r] \arrow[d] & Z^{n-1} \arrow[r] \arrow[d] & 0 \\
		0 \arrow[r] & X^n \arrow[r] & Y^n \arrow[r] & Z^n \arrow[r] & 0
	\end{tikzcd}\]%
	\footnote{Translator's note (O014-C043): the source states the sequence
	$\Coker(d_X^n)\to\Coker(d_Y^n)\to\Coker(d_Z^n)\to0$, while the vertical
	arrows in the diagram immediately following it are $d_X^{n-1}$,
	$d_Y^{n-1}$, and $d_Z^{n-1}$. This translation aligns the cokernel indices
	with the diagram; renaming the index then gives the row indexed by $n-2$
	used below.}
	Similarly, there is an exact sequence
	$0\to\Ker(d_X^n)\to\Ker(d_Y^n)\to\Ker(d_Z^n)$. The two sequences fit
	into the following commutative diagram with exact rows:
% segment-id: o014.li2.chapter3.abelian-category-complexes.g006
	\begin{equation}\label{eqn:long-exact-sequence-ses-aux}\begin{tikzcd}
		& \Coker(d_X^{n-2}) \arrow[r] \arrow[d] & \Coker(d_Y^{n-2}) \arrow[r] \arrow[d] & \Coker(d_Z^{n-2}) \arrow[d] \arrow[r] & 0 \\
		0 \arrow[r] & \Ker(d_X^n) \arrow[r] & \Ker(d_Y^n) \arrow[r] & \Ker(d_Z^n) &
	\end{tikzcd}\end{equation}
	The vertical arrows are induced respectively by $d_X^{n-1}$,
	$d_Y^{n-1}$, and $d_Z^{n-1}$ and have epi--mono factorizations of the
	form
	$\Coker(d^{n-2})\stackrel{d^{n-1}}{\twoheadrightarrow}
	\Image(d^{n-1})\hookrightarrow\Ker(d^n)$, with subscripts suppressed.

	By Lemma~\ref{prop:mono-epi-ker-coker} (iii), the epimorphism
	$\Coker(d^{n-2})\twoheadrightarrow\Image(d^{n-1})$ determines the
	kernels of the vertical arrows in
	\eqref{eqn:long-exact-sequence-ses-aux}; then
	\eqref{eqn:homology-3} gives in turn $\Hm^{n-1}(X)$, $\Hm^{n-1}(Y)$,
	and $\Hm^{n-1}(Z)$. Likewise, the monomorphism
	$\Image(d^{n-1})\hookrightarrow\Ker(d^n)$ determines their cokernels,
	giving in turn $\Hm^n(X)$, $\Hm^n(Y)$, and $\Hm^n(Z)$. Applying
	Theorem~\ref{prop:snake-lemma} now yields the connecting morphism
	$\delta^{n-1}:\Hm^{n-1}(Z)\to\Hm^n(X)$ and the portion of the long exact
	sequence containing the terms $\Hm^n$ and $\Hm^{n-1}$. Its functoriality
	follows from Remark~\ref{rem:connecting-canonical}.
\end{proof}

% segment-id: o014.li2.chapter3.abelian-category-complexes.d001
\begin{definition}[Quasi-isomorphism]\label{def:quasi-isomorphism}
	\index{quasi-isomorphism}
	Let $f:X\to Y$ be a morphism in $\cate{C}(\mathcal{A})$. If $\Hm^n(f)$
	is an isomorphism for every $n\in\Z$, then $f$ is called a
	\emph{quasi-isomorphism}.
\end{definition}

% segment-id: o014.li2.chapter3.abelian-category-complexes.pr004
\begin{proposition}\label{prop:null-homotopic-cohomology}
	Let $f:X\to Y$ be a morphism in $\cate{C}(\mathcal{A})$ and let
	$n\in\Z$. If $f$ is null-homotopic, then $\Hm^n(f)=0$.
\end{proposition}

% segment-id: o014.li2.chapter3.abelian-category-complexes.pf004
\begin{proof}
	Choose $h\in\Hom^{-1}(X,Y)$ such that $f=d_Yh+hd_X$. The composite
	$\Ker(d_X^n)\hookrightarrow X^n\xrightarrow{h^{n+1}d_X^n}Y^n$ is zero.
	On the other hand, $d_Y^{n-1}h^n$ factors through
	$\Image(d_Y^{n-1})$. Together these facts give $\Hm^n(f)=0$.
\end{proof}

% segment-id: o014.li2.chapter3.abelian-category-complexes.co001
\begin{corollary}
	For every $n\in\Z$, the cohomology functor
	$\Hm^n:\cate{C}(\mathcal{A})\to\mathcal{A}$ factors uniquely through
	$\cate{K}(\mathcal{A})$. The notion of quasi-isomorphism therefore extends
	to morphisms in $\cate{K}(\mathcal{A})$. If $f:X\to Y$ is an isomorphism
	in $\cate{K}(\mathcal{A})$, then $f$ is a quasi-isomorphism.
\end{corollary}

% segment-id: o014.li2.chapter3.abelian-category-complexes.pf005
\begin{proof}
	Combine Propositions~\ref{prop:homotopy-category-universal-property} and
	\ref{prop:null-homotopic-cohomology}.
\end{proof}
